\section{Introduction}

A neural intrusion detection system (IDS) can perform well on held-out traffic
without establishing that its feature attributions are reliable evidence for
auditing individual decisions. An analyst inspecting a flagged flow needs to
know what an explanation measures, whether it can be reproduced, and whether
the attributed features affect the model under an explicit intervention.
Trust motivates these questions; neither predictive accuracy nor the presence
of an explanation answers them. Benchmark performance also need not transfer
to the conditions in which an IDS is deployed~\cite{sommer2010outside}.

Post-hoc methods produce different accounts of a fitted model. A local
surrogate coefficient, a path-integrated gradient and a replacement effect
are not interchangeable measurements. Repeating an account can establish
reproducibility without establishing that its ordering predicts a functional
response. Conversely, a deletion experiment depends on what replaces the
deleted coordinates~\cite{rong2022road,wang2024deletion}. An audit must therefore
separate properties of the explanation from choices made by its evaluator.

We ask: \emph{To what extent do post-hoc explanations provide reliable evidence
for auditing neural intrusion detection decisions?} We operationalise reliability
through three questions, not a composite trust score:
\textbf{RQ1}: how reproducible are attributions under repeated explanation
sampling or changes in numerical integration, for a fixed model and target?
\textbf{RQ2}: how do attribution-guided interventions compare with matched
random controls, and how does this response depend on the replacement operator?
\textbf{RQ3}: how does the response vary across correct and incorrect decisions,
including false positives and false negatives?

The study reorganises existing experiments on two binary traffic benchmarks,
two neural architectures, five training seeds and five explainers. Its
contribution is a condition-resolved audit: metric-specific coverage accompanies
reproducibility; paired functional effects are evaluated under three distinct
operators; and an exploratory outcome analysis states its restricted method
coverage and sample support. We evaluate evidence about frozen models, not
analyst trust, incident reduction, or a criterion for declaring an IDS safe.
